
%%FILE blueprint/src/content.tex
% In this file you should put the actual content of the blueprint.
% It will be used both by the web and the print version.
% It should *not* include the \begin{document}
%
% If you want to split the blueprint content into several files then
% the current file can be a simple sequence of \input. Otherwise It
% can start with a \section or \chapter for instance.
%

\chapter{Formalization Methodologies}
\input{modules/ComputationalRel}

\chapter{Midgard Modules Formalization}
\input{modules/Cardano}
\input{modules/List}
\input{modules/Parameters}
\input{modules/OperatorDirectory}
\input{modules/Scheduler}
\input{modules/StateQueue}
\input{modules/ProofProtocol}
\input{modules/Settlement}
\input{modules/UserEvent}

\chapter{Midgard Composition}

The Midgard protocol emerges from the complex interaction between the modules
defined in the previous chapter.
%
Here we list the interactions implemented using the combinators described in
Section~\ref{sec:computationalrelations}.
%
Mainly sub-interface restrictions and tensor products of machines and relations.
%
We use parallel composition of relation explicitly using the disjoint unions of
actions.

For each module simultaneously interacting with another, we need to duplicate actions redeemers
(arguments) and explicitly state extra conditions.
%
Moreover, we need to also restrict the module to the operations it can perform
alone, without interacting with other modules.
%
This is a consequence of extra-modularization, however we opted to follow the
specification as closely as possible.
%
An advantage is that we can define each relation and their function computing them in isolation,
and then, we know that the composition and restriction of computable inductive
relations is also a computable inductive relation.
%
Finally, we can use the Blaster tool to test properties over the composition of
different modules out-of-the-box.

The implementation follows the same structure as the other modules.
%
We refer the reader to see the following files:
\begin{description}
\item[Composed Midgard State] ``FMMidgard/State.lean''
\item[Composed Midgard Relation] ``FMMidgard/Relation.lean''
\item[Executing Function]``FMMidgard/StateMachine.lean''
  \end{description}

\chapter{Extras}
\section{Issues}
\label{sec:Issues}

Small section listing the submitted issues.

\begin{description}
  \item[Registered Operators Node default values]\url{https://github.com/Anastasia-Labs/midgard/issues/295}
  \item[Init Scheduler Datum]\url{https://github.com/Anastasia-Labs/midgard/issues/296}
  \item[Operator Inactivity Timeout]\url{https://github.com/Anastasia-Labs/midgard/issues/297}
  \item[Registered Operators Eligible for Activation Criteria and Scheduler Rewind]\url{https://github.com/Anastasia-Labs/midgard/issues/311}
  \item[State Queue Root Node Initialization Values]\url{https://github.com/Anastasia-Labs/midgard/issues/313}
  \item[State Queue Root node unspecified values when merging confirmed blocks]\url{https://github.com/Anastasia-Labs/midgard/issues/314}
  \item[Breaking Registered Operators Directory by registering twice]\url{https://github.com/Anastasia-Labs/midgard/issues/318}
  \item[Scheduler gets stuck if last active operator is operating and retires]\url{https://github.com/Anastasia-Labs/midgard/issues/391}
  \item[Anyone can retire active operators]\url{https://github.com/Anastasia-Labs/midgard/issues/395}
  \item[Anyone can claim bonds of retired operators.]\url{https://github.com/Anastasia-Labs/midgard/issues/397}
  \item[Check that \(\mathsf{start\_time} < \mathsf{end\_time}\) in block headers]\url{https://github.com/Anastasia-Labs/midgard/issues/400}
  \item[Update the spec specifying the validity range limit parameter]\url{https://github.com/Anastasia-Labs/midgard/issues/401}
  \item[Escape Hatch]\url{https://github.com/Anastasia-Labs/midgard/issues/402}
  \item[Rollbacks and Staking Script Credentials]\url{https://github.com/Anastasia-Labs/midgard/issues/403}
  \item[Non-refundable Deposits with Refund Address and Datum fields]\url{https://github.com/Anastasia-Labs/midgard/issues/405}
   \end{description}


%%FILE blueprint/src/modules/Cardano.tex
\section{Cardano}

In this section, we describe how we model Cardano and what assumptions we place
on the main layer.
%
All these constructs can be found inside the ``FMMidgard/Cardano'' folder.

\subsection{Basic Cardano Types}

We model some Cardano types using Lean4 types.
%
\begin{itemize}
    \item Cardano \emph{addresses} simply as Lean4 \lean{String}
    \item Cardano value (ADA and others) as Lean4 \lean{Nat}
    \item Cardano data (bytestrings) as Lean4 \lean{String}
    \item Hashes simple as Lean4 \lean{Nat}
\end{itemize}
%
The main reason to keep so simple types is to keep the model simple and
compatible with the Blaster tool.

In the real implementation and more concrete models or formalizations of the
protocol should implement \emph{finite} types.
%
Hashes, bytestrings, and addresses in Cardano have a fixed length, and we are
not taking that into account in this formalization.

\paragraph*{Transactions}

Midgard is a timed protocol and it is in sync with Cardano.
%
All transactions in Midgard use \emph{transaction validity intervals} to
measure time.
%
Cardano guarantees that time progresses and that the lower bound is less than
current time and upper bound is more than current time.

Midgard relies on Cardano for authentication of transactions, and thus, it
checks who is signing a given transaction.

For simplicity, we have two kind of transactions.
%
One basic transaction exposing just time intervals and who is signing a given
transaction, for simplicity we assume there is one key.

\begin{definition}{Transaction Information}
\label{dfn:InputInfo}
Transaction are a tuple providing:
\begin{description}
  \item[validity internal]
        execution time bounds. When missing, they are assumed to by minus
        infinity and plus infinity.
  \item[value] the amount of value in the transaction
  \item[fee] the fee of the transaction
  \item[validator] who is signing the transaction
  \item[transaction id]
  \end{description}
\end{definition}

There is a complex definition of transaction, closer to the real implementation.
%
However, we only use this complex definition to mechanize the \emph{User Event}
module.

\subsection{Tokens}

Cardano tokens abstraction.
%
Computations are token threaded, in the sense that since we live in the UTxO
model, we pass tokens around to indicate and authenticate computations and states.
%
Midgard is not an exception, so we added a mechanization of tokens.

\martin{I skipped lean \emph{Strings} because the Blaster tool does not support
them, but as soon as we can use, we need them for token names.}

\begin{definition}[Tokens]
  Tokens are structures containing a name, (polymorphic) data and id.
  \label{dfn:Tokens.Elem.TokenD}
  \lean{Tokens.Elem.TokenD}
  \leanok
  Tokens are records with:
  \begin{description}
    \item[id] same as its reference
    \item[data] keeping internal data
    \item[name] the name of the token
  \end{description}
  \end{definition}
%
The id here is a mechanization artifact, in Cardano we do not need this as there
are other mechanisms. 

When minting new tokens, a new id is generated and guaranteed to be unique.
%
The action of \emph{having/presenting} a token here is formalized has
\emph{knowing} the correct id and we delegate the ownership model to lower-level
formal models.
%
In Cardano, this is not necessary since the UTxO-extended model already provides it.

Different implementations are possible, in this model we use lists containing
elements that may not be there (i.e. they were burnt.)
%
The element position in the list is its id, and thus, all elements have a unique
id since we only append elements.
%
When minting, we just append a new value, and when burning, we just mark it as
no element.
%
Therefore, we do not reuse ids, and the implementation is simply enough that the
Blaster tool can be used to reason about modules using this structure.



Properties about this module are needed to prove other properties of other
modules.
%
Instead of listing them and link them to the document, we just refer to them
using the datatype.

\martin{This is instrumental to the library, although it may be good to notice
what properties about tokens we need to prove properties about Midgrad.}

\section{Hashes}
%
Midgard makes heavy use of hashes and other data structures.
%
As part of Cardano we assume there are some hashes employed as keys and identifiers.
%
Again, we abstract those types as Lean4 \lean{Nat} for now, because we do not
have a good way of representing Hash functions and we want to be compatible with
the Blaster tool.


%%FILE blueprint/src/modules/ComputationalRel.tex
\section{Methodology}

In this section, we describe the methodology we use to define the different
Midgard Modules.
%
We follow the approach described by other members of the formal methods team:
\emph{computable inductive relations.}
%
We define an inductive \emph{partial} relation describing the different
operations modules have, and then, a computable function realizing
such \emph{partial} relation.
%
Functions realizing partial functions are \emph{partial functions}.
%
If the values are related, the function successfully shows that, if not, it
returns an error.

Each Midgard module defines:
%
\begin{description}
    \item[Inductive relation:] a Lean4 inductive relation exposing each module
        operations (minting policies and spending validators).
    \item[Computational function:] a Lean4 computational function following
        closely what the inductive relation does.
    \item[Properties:] a set of properties, beginning by showing that the
        computable function realizes the inductive relation of the module.
\end{description}

Significant properties are usually found by a rigorous \emph{Safety Analysis.}
%
Such analysis is left outside this document.

\subsection{Computational Relations.}\label{sec:computationalrelations}

Inductive Relations and computable functions have advantages.
%
On the one hand, defining inductive relations is the standard way to reason about the evolution
of programs in the dependent type languages community (Agda/Rocq/etc.)
%
While, on the other hand, to reason about reactive programs, it is very useful
to have actual functions to execute.
%
For example, we can use \emph{Blaster} to find counter-examples and traces
leading to invalid states.

Following the Agda-Ledger approach~\cite{Ledger.FMBC2024}, we define a
class stating that a relation is realizable by a function.

\begin{definition}[Computational Relation]
  A relation \(\mathcal{R} \subseteq E \times A \times S \times S \) is
  computational if and only if there is a function
  \(f : E \times A \times S \to S_{\mathsf{Errors}}\) such that
  \[
    \forall (e \in E) (a \in A) (s_{i}, s_{o} \in S),
    (e, a, s_{i}, s_{o}) \in \mathcal{R}  \iff
    f(e,a,s_{i}) = \eta(s_{o})
  \]
  \label{dfn:computational_relation}
  \lean{ComputationalRel}
  \leanok
\end{definition}

Since we want to predicate over reachable states, the ones following the
operations described in a module, we define them as just the ones after
executing some steps.

\begin{definition}[From Initial State]
  \lean{StepMachines.StepMachine.FromInit}
  \leanok
  \label{dfn:ComputationalRelation.StateMachine.FromInit}
  %%
  Given a State Machine \(\mathsf{sm}\), we say that a state \(s\) is reach from
the initial state \(\mathsf{sm.init}\) if there are 0 or more \(\mathsf{sm}\)-steps from
\(\mathsf{sm.init}\) to \(s\).
  \end{definition}

Some properties are not valid for the initial state but for the rest of the
reachable states.

\begin{definition}[After Initial State]
  \lean{StepMachines.StepMachine.AfterInit}
  \leanok
  \label{dfn:ComputationalRelation.StateMachine.AfterInit}
  %%
  Given a State Machine \(\mathsf{sm}\), we say that a state \(s\) is reach from
the initial state \(\mathsf{sm.init}\) if there are 1 or more \(\mathsf{sm}\)-steps from
\(\mathsf{sm.init}\) to \(s\).
  \end{definition}


\paragraph{Induction}
%
We can prove invariants over these relation by structural induction over each
possible action.

\begin{definition}[Induction Step Machine]
  \lean{StepMachines.ApplyInd}
  \leanok
  \label{dfn:Relations.StepMachines.ApplyInd}
  Given a predicate that is valid for the initial state and as an invariant over
the relation steps, we can prove the predicate holds for all reachable states.
  \end{definition}

We can prove the same for states after the initial one.

\begin{definition}[Induction After Initial State Machine]
  \lean{StepMachines.ApplyFromInit}
  \leanok
  \label{dfn:Relations.StepMachines.ApplyFromInit}
  \end{definition}


\paragraph{Composition}

To define the final Midgard computable inductive relation, we need to compose
the different sub-modules.

We define the parallel and product composition of machines, plus a conditional
restriction operation.
%
The restriction operation enable us to define simultaneous operations adding
constraints related to sub-machines.

\begin{definition}[Relation Tensor Product]
  \lean{Products.Joint}
  \leanok
  \label{dfn:Relations.Products.Joint}
  \end{definition}

\begin{definition}[Relation Conditional Product]
  \lean{Products.Conditional_Joint}
  \leanok
  \label{dfn:Relations.Products.Conditional_Joint}
  \end{definition}

In practice, it is also useful to have interface restriction/conditional
restriction over relations.
%
This way, instead of using conditional products, we can product and then restrict.

\begin{definition}[Relational Interface Restriction]
  \lean{InterfaceRestrictions.ConditionalSub}
  \leanok
  \label{dfn:Relations.InterfaceRestrictions.ConditionalSub}
  \end{definition}

On the computable functions side, we replicate these definitions.

\begin{definition}[Machine Product]
  \lean{CompositionAndOps.product_comp}
  \leanok
  \label{dfn:ComputationalRelation.CompositionAndOps.product_comp}
  \end{definition}

We have two parallel composition of machines, one assuming there is no sharing
between the two machines and another where both machines can see each other
state.

\begin{definition}[Machine Parallel Composition]
  \lean{CompositionAndOps.parallel_comp}
  \leanok
  \label{dfn:ComputationalRelation.CompositionAndOps.parallel_comp}
  \end{definition}

\begin{definition}[Machine Parallel Composition Sharing State]
  \lean{CompositionAndOps.parallel_comp_sharing_state}
  \leanok
  \label{dfn:ComputationalRelation.CompositionAndOps.parallel_comp_sharing_state}
  \end{definition}

Finally, we also have conditional interface restriction at the machine level.

\begin{definition}[Conditional Interface Restriction]
  \lean{CompositionAndOps.conditional_sub}
  \leanok
  \label{dfn:ComputationalRelation.CompositionAndOps.conditional_sub}
  \end{definition}

All the above operations define computable inductive relations when applied to such relations.

\begin{definition}[Computational Product Composition]
  \lean{CompositionAndOps.computational_product}
  \leanok
  \label{dfn:ComputationalRelation.CompositionAndOps.computational_product}
  \uses{dfn:ComputationalRelation.CompositionAndOps.product_comp,dfn:Relations.Products.Joint, dfn:computational_relation}
  \end{definition}

\begin{definition}[Computational Parallel Composition]
  \lean{CompositionAndOps.computational_parallel}
  \leanok
  \label{dfn:ComputationalRelation.CompositionAndOps.computational_parallel}
  \uses{dfn:ComputationalRelation.CompositionAndOps.parallel_comp, dfn:computational_relation}
  \end{definition}


\begin{definition}[Computational Conditional Subinterfacing]
  \lean{CompositionAndOps.computational_conditional}
  \leanok
  \label{dfn:ComputationalRelation.CompositionAndOps.computational_conditional}
  \uses{dfn:ComputationalRelation.CompositionAndOps.conditional_sub,dfn:Relations.InterfaceRestrictions.ConditionalSub, dfn:computational_relation}
  \end{definition}


%%FILE blueprint/src/modules/List.tex
\section{Lists}

The Midgard protocol uses singly linked lists as the main data structure for the
operator directory and state queue.
%
The specification defines two types of lists: key-unordered and key-ordered lists.
%
The key-unordered list data structure \emph{assumes} keys are unique, while the
key-ordered list data structure inserts them in strict order guaranteeing they
are unique and sorted.
%
Each element in the key-ordered list represents the element and a range of
elements that are not in the list.

The specification designers noticed that when the key-uniqueness
assumption is violated, the data structure diverges from usual list
implementations.
%
In particular, traversals are no deterministic.
%
At first, we used \emph{Lean4} lists and abstracted them away.
%
However, the Operator Directory module let operators to register more than once
making the acceptance of duplicates a desired (but penalized) behaviour.
%
When keys (operators identifier keys) can register more than once, they break
one of the main assumptions of the data structure.
%
Moreover, it is not just that traversals are not deterministic, it is possible
to \textbf{break} the structure.

To detect violations, we employ the Plausible framework, a simple
property-based framework in Lean4, for more information see
Section~\ref{sec:list:plausible}.
%
As result, we can test for different properties generating counter-examples
(traces in this case) leading to errors in Midgard.
%
We documented this pattern and submitted the issues found (see \ref{sec:Issues}.)

% Since we though the key unordered lists specification was fishy, we implemented
% them in more detail.
% %
% After implementing them, we can \emph{generate} test cases and break the
% list structure.
% %
% and we could generate another run at the Operator Directory
% level too (See sections \ref{sec:list:plausible, sec:Issues}.)


To add elements in either of these lists, we need to mint a token with the
information we want to add, and to remove an element, we simply burn such token.
%
To prove that an element is in the list, we just show a token that contains its 
as the mere existence of the token shows membership plus the correct
\emph{Spending/Minting Policies}.
%
To prove that an element is not in the list, we need to show the element is not
where it should be, this requires the list to be sorted.
%
Membership is possible on sorted/unsorted lists, while non-membership requires
the list to be sorted.

\subsection{Plots}

\paragraph*{Input-Outputs.}
Consuming utxos and data for each policies.

\includegraphics[width=\textwidth]{io.linkedlist}


\subsection{Key Unordered Lists}

To simulate tokens in our formalization, we simply use a (finite) map from
identifiers to \emph{tokens.}
%
These tokens have internal data, and in the case of lists are structures defined
as follows:
%
\begin{definition}[List Element]
  \label{dfn:MUList.State.Elemt}
  \lean{MUList.State.Elemt}
  \leanok
  List nodes are records with:
  \begin{description}
    \item[key] (optional) indicates the current node key,
    \item[link] (optional) indicates the next key in the list,
    \item[data] polymorphic element depending on the use of the list.
  \end{description}
  \end{definition}

\begin{definition}[Onchain Lists]
  \label{dfn:MUList.State.ListToken}
  \lean{MUList.State.ListToken}
  \leanok
  \uses{dfn:Tokens.Elem.TokenD, dfn:MUList.State.Elemt}
  Lists are defined as a collection of \(\textsf{Tokens}\)
  with internal data \(\textsf{Elemt}\)

  Internally, these collections are mechanized as maps from \(\textsf{TId}\) to
\(\textsf{Tokens}\) or burnt tokens.
  \end{definition}

\subsubsection{State Machine}

We define the following transitions (or actions), we have over these lists.

\includegraphics[width=\textwidth]{linked.list}

\begin{definition}[Actions]
  \label{dfn:MUList.State.Actions}
  \lean{MUList.State.Actions}
  \leanok
  We have the following actions.
  \begin{description}
    \item[Modify Data] modifies the data field of a node in the list
    \item[Unsafe Append] adds an element at the end of the list
    \item[Unsafe Prepend] adds an element at the beginning of the list
    \item[Remove] removes an element from the list
    \end{description}
  \end{definition}

Each actions have different arguments (redeemers).
%
Since we do not traverse the list onchain (traversals are not deterministic),
for each action we need to provide the tokens to be updated.
%
For example, to add an element to the end of the list,
\(\mathsf{unsafe append}\), we need to provide the last element of the list.
%
In the formalization, providing elements is modelled as giving the
\(\mathsf{TId}\) of the token, while in Cardano, we just provide it.

We defined the inductive relation.

\begin{definition}[Unordered List Relation]
  \label{dfn:MUList.Relation.Rel}
  \lean{MUList.Relation.Rel}
  \uses{dfn:MUList.State.ListToken, dfn:MUList.State.Actions}
  \leanok
  \end{definition}

\begin{definition}[Unordered List computation]
  \label{dfn:MUList.Computation.next}
  \lean{MUList.Computation.next}
  \uses{dfn:MUList.State.ListToken, dfn:MUList.State.Actions}
  \leanok
  \end{definition}

%
\begin{theorem}[Computational Relation]
  \label{thm:MUList.Computational.ComputationalRelation}
  \lean{MUList.Computational.ComputationalRelation}
  \uses{dfn:MUList.Computation.next, dfn:MUList.Relation.Rel}
  \leanok
  \end{theorem}


\subsubsection*{Consistency Proofs}

\begin{theorem}[Lists are never empty]
  \label{thm:MUList.Properties.NotEmpty}
  \lean{MUList.Properties.NotEmpty}
  \leanok
  There is always a token in the list.
  \end{theorem}

\begin{proof}
  \proves{thm:MUList.Properties.NotEmpty}
  The root node is always alive, or other nodes are added.
  \end{proof}

  
\begin{theorem}
  \label{thm:MUList.Properties.not_remove_root}
  \lean{MUList.Properties.not_remove_root}
  \leanok
  Root nodes cannot be removed.
  \end{theorem}

\begin{theorem}
  \label{thm:MUList.Properties.one_root}
  \lean{MUList.Properties.one_root}
  \leanok
  Root node remains the same.
\end{theorem}


We prove some correctness properties regarding the use of tokens.

\begin{theorem}[Reference and Id equiv]
  \label{thm:MUList.Properties.proof_id_match}
  \lean{MUList.Properties.proof_id_match}
  \leanok
  The reference assigned when adding elements and the id field in the token matches.
  \end{theorem}


\paragraph*{Prepend Properties}

\begin{theorem}[Enabled Prepend]
  \label{thm:MUList.Properties.EnabledPrepend}
  \lean{MUList.Properties.EnabledPrepend}
  \leanok
  %%
  Prepending elements to lists is always enabled if we give the root node.
  %%
  \end{theorem}

\begin{theorem}[Prepend Prepends]
  \label{thm:MUList.Properties.Prepend_Props}
  \lean{MUList.Properties.Prepend_Props}
  \uses{dfn:MUList.State.Actions,dfn:MUList.Computation.next}
  \leanok
  After prepending an element, we have a token showing its membership.
  %
  Furthermore, the root node is a valid anchor for the recently prepended node.
  \end{theorem}

\begin{theorem}[Prepend Prepends Computation]
  \label{thm:MUList.Properties.Prepend_prepends'}
  \lean{MUList.Properties.Prepend_prepends'}
  \uses{dfn:MUList.State.Actions,dfn:MUList.Computation.next}
  \leanok
  Same as before but mapped to list functions: \(\mathsf{is\_member,is\_anchor,is\_root\_node}\).
  \end{theorem}

\begin{theorem}[Prepend Prepends Relation]
  \label{thm:MUList.Properties.Prepend_prepends}
  \lean{MUList.Properties.Prepend_prepends}
  \uses{dfn:MUList.State.Actions,dfn:MUList.Relation.Rel}
  \leanok
  Same as before using the relation. 
  \end{theorem}

\begin{theorem}[Prepend New Element]
  \label{thm:MUList.Properties.Prepend_prepends_new}
  \lean{MUList.Properties.Prepend_prepends_new}
  \uses{dfn:MUList.State.Actions,dfn:MUList.Relation.Rel}
  \leanok
  The element added is new.
  \end{theorem}

\begin{theorem}[Prepend Others ]
  \label{thm:MUList.Properties.Prepend_others_same}
  \lean{MUList.Properties.Prepend_others_same}
  \uses{dfn:MUList.State.Actions,dfn:MUList.Relation.Rel}
  \leanok
  After prepends all other elements, neither the root nor the just added, remain the same.
  \end{theorem}

\paragraph*{Append Properties}

\begin{theorem}[Enabled Append]
  \label{thm:MUList.Properties.EnabledAppend }
  \lean{MUList.Properties.EnabledAppend }
  \uses{dfn:MUList.State.Actions,dfn:MUList.Computation.next}
  \leanok
  \end{theorem} 


\begin{theorem}[Append After state Properties]
  \label{thm:MUList.Properties.Append_Props}
  \lean{MUList.Properties.Append_Props}
  \uses{dfn:MUList.State.Actions,dfn:MUList.Computation.next}
  \leanok
  \end{theorem} 

\begin{theorem}[Append After State Properties (list operations)]
  \label{thm:MUList.Properties.Append_appends}
  \lean{MUList.Properties.Append_appends}
  \uses{dfn:MUList.State.Actions,dfn:MUList.Relation.Rel}
  \leanok
  \end{theorem} 

\paragraph*{Removes Properties}

\begin{theorem}[Enabled Remove]
  \label{thm:MUList.Properties.EnabledRemove }
  \lean{MUList.Properties.EnabledRemove }
  \uses{dfn:MUList.State.Actions,dfn:MUList.Computation.next}
  \leanok
  \end{theorem} 
\begin{theorem}[Enabled Remove Relation]
  \label{thm:MUList.Properties.EnabledRemove'}
  \lean{MUList.Properties.EnabledRemove'}
  \uses{dfn:MUList.State.Actions,dfn:MUList.Relation.Rel}
  \leanok
  \end{theorem} 

\begin{theorem}[Remove Member]
  \label{thm:MUList.Properties.Removes_members}
  \lean{MUList.Properties.Removes_members}
  \uses{dfn:MUList.State.Actions,dfn:MUList.Relation.Rel}
  \leanok
  After removing an element, the element must have had been a member of the
list.
  \end{theorem} 
  
\begin{theorem}[Remove Anchor]
  \label{thm:MUList.Properties.Removes_anch}
  \lean{MUList.Properties.Removes_anch}
  \uses{dfn:MUList.State.Actions,dfn:MUList.Relation.Rel}
  \leanok
  If we can remove an element, then the anchor proposed is a valid anchor of the
element to be removed.
  \end{theorem} 

\begin{theorem}[Remove Not Self (function)]
  \label{thm:MUList.Properties.NotRemoveself'}
  \lean{MUList.Properties.NotRemoveself'}
  \uses{dfn:MUList.State.Actions,dfn:MUList.Computation.next}
  \leanok
  We cannot remove an element using itself as an anchor.
  \end{theorem} 

\begin{theorem}[Remove Not Self (Relation)]
  \label{thm:MUList.Properties.NotRemoveSelf}
  \lean{MUList.Properties.NotRemoveSelf}
  \uses{dfn:MUList.State.Actions,dfn:MUList.Relation.Rel}
  \leanok
  We cannot remove an element using itself as an anchor.
  \end{theorem} 
  
\begin{theorem}[Removed element is not a member.]
  \label{thm:MUList.Properties.Removes_finds_none}
  \lean{MUList.Properties.Removes_finds_none}
  \uses{dfn:MUList.State.Actions,dfn:MUList.Relation.Rel}
  \leanok
  After removing an element, the element is not a member anymore.
  \end{theorem} 

\begin{theorem}
  \label{thm:MUList.Properties.Removes_others}
  \lean{MUList.Properties.Removes_others}
  \uses{dfn:MUList.State.Actions,dfn:MUList.Relation.Rel}
  \leanok
  After removing an element, all other elements remain the same, but the removed element and its anchor.
  \end{theorem} 
  
\begin{theorem}
  \label{thm:MUList.Properties.Removes_update_anch_mem_and_key}
  \lean{MUList.Properties.Removes_update_anch_mem_and_key}
  \uses{dfn:MUList.State.Actions,dfn:MUList.Relation.Rel}
  \leanok
  After removing an element, its anchor links to the element the removed node
was linking to, and everything else remains the same.
  \end{theorem} 

\begin{theorem}
  \label{thm:MUList.Properties.Removes_update_anch_mem_and_key'}
  \lean{MUList.Properties.Removes_update_anch_mem_and_key'}
  \uses{dfn:MUList.State.Actions,dfn:MUList.Relation.Rel}
  \leanok
  Same as before but in a more functional way.
  \end{theorem} 

\begin{theorem}
  \label{thm:MUList.Properties.Removes_update_anch_mem}
  \lean{MUList.Properties.Removes_update_anch_mem}
  \uses{dfn:MUList.State.Actions,dfn:MUList.Relation.Rel}
  \leanok
  Anchors membership remains the same.
  \end{theorem} 

\begin{theorem}
  \label{thm:MUList.Properties.Removes_same_anch_key}
  \lean{MUList.Properties.Removes_same_anch_key}
  \uses{dfn:MUList.State.Actions,dfn:MUList.Relation.Rel}
  \leanok
  Anchors key remains the same.
  \end{theorem} 

\begin{theorem}
  \label{thm:MUList.Properties.Removes_others_keys}
  \lean{MUList.Properties.Removes_others_keys}
  \uses{dfn:MUList.State.Actions,dfn:MUList.Relation.Rel}
  \leanok
  Elements aside from the anchor and removed tokens keep the same key.
  \end{theorem} 

\begin{theorem}
  \label{thm:MUList.Properties.Removes_others_mem}
  \lean{MUList.Properties.Removes_others_mem}
  \uses{dfn:MUList.State.Actions,dfn:MUList.Relation.Rel}
  \leanok
  Elements aside from the anchor and removed tokens keep the same.
  \end{theorem} 

\begin{theorem}
  \label{thm:MUList.Properties.Removes_members_is_mem}
  \lean{MUList.Properties.Removes_members_is_mem}
  \uses{dfn:MUList.State.Actions,dfn:MUList.Relation.Rel}
  \leanok
  If an element is a member of the resulting list after removing an element,
then it is a member of the original list.
  \end{theorem} 

\begin{theorem}
  \label{thm:MUList.Properties.Removes_members_anch}
  \lean{MUList.Properties.Removes_members_anch}
  \uses{dfn:MUList.State.Actions,dfn:MUList.Relation.Rel}
  \leanok
  After removing, the anchor links gets updated properly.
  \end{theorem} 

\begin{theorem}
  \label{thm:MUList.Properties.Remove_removes}
  \lean{MUList.Properties.Remove_removes}
  \uses{dfn:MUList.State.Actions,dfn:MUList.Relation.Rel}
  \leanok
  After removing, the removed element is not a member anymore.
  \end{theorem} 

\begin{theorem}
  \label{thm:MUList.Properties.Removes_get_not_removed}
  \lean{MUList.Properties.Removes_get_not_removed}
  \uses{dfn:MUList.State.Actions,dfn:MUList.Relation.Rel}
  \leanok
  If an element \(e\) is a member of a list after removing an element \(r\), then the
element \(e\) is not the same as the removed one, i.e. \(e \neq r\).
  \end{theorem} 

%%
\paragraph*{Modify Data Properties}

\begin{theorem}
  \label{thm:MUList.Properties.ModifyData_same_others}
  \lean{MUList.Properties.ModifyData_same_others}
  \uses{dfn:MUList.State.Actions,dfn:MUList.Relation.Rel}
  \leanok
  After modifying an element, all the rest remain the same.
  \end{theorem} 

\begin{theorem}
  \label{thm:MUList.Properties.ModifyData_find_token}
  \lean{MUList.Properties.ModifyData_find_token}
  \uses{dfn:MUList.State.Actions,dfn:MUList.Relation.Rel}
  \leanok
  If we can modify an element, the element is a member of the list before and
after modifying it.
  \end{theorem} 

\begin{theorem}
  \label{thm:MUList.Properties.ModifyData_modified}
  \lean{MUList.Properties.ModifyData_modified}
  \uses{dfn:MUList.State.Actions,dfn:MUList.Relation.Rel}
  \leanok
  Modifying the element modifies the data of the element as expected.
  \end{theorem} 

\begin{theorem}
  \label{thm:MUList.Properties.ModifyData_modified_stronger}
  \lean{MUList.Properties.ModifyData_modified_stronger}
  \uses{dfn:MUList.State.Actions,dfn:MUList.Relation.Rel}
  \leanok
  Stronger version of previous theorem, stating that everything else remains the same.
  \end{theorem} 
\begin{theorem}
  \label{thm:MUList.Properties.ModifyData_modified_stronger'}
  \lean{MUList.Properties.ModifyData_modified_stronger'}
  \uses{dfn:MUList.State.Actions,dfn:MUList.Relation.Rel}
  \leanok
  Similar to the previous theorem.
  \end{theorem} 

\begin{theorem}
  \label{thm:MUList.Properties.ModifyData_same_members}
  \lean{MUList.Properties.ModifyData_same_members}
  \uses{dfn:MUList.State.Actions,dfn:MUList.Relation.Rel}
  \leanok
  After modifying an element, all other elements are still in the list.
  \end{theorem} 

\subsection{Key Ordered Lists }

Internally, the Midgard Ordered lists are the same as the Unoredred ones, where
the main difference is in the set of actions and set of validators.


\begin{definition}[Actions]
  \label{dfn:MOList.State.Actions}
  \lean{MOList.State.Actions}
  \leanok
  We have the following actions.
  \begin{description}
    \item[Modify Data] modifies the data field of a node in the list
    \item[Safe Append] adds an element at the end of the list
    \item[Safe Prepend] adds an element at the beginning of the list
    \item[Insert] insert an element after a given anchor node
    \item[Remove] removes an element from the list
    \end{description}
  \end{definition}

Since we have a new operation, technically we cannot map properties of
Unoredered lists as easily.
%
We would need to define some notion of simulation/bisimulation or other notions,
stating that all reachable states on both structures are reachable on the other.
%
However, this is not strictly necessary to prove properties about them.


\begin{definition}[Inductive Relation]
  \label{dfn:MOList.Relation.Rel}
  \lean{MOList.Relation.Rel}
  \leanok
  Inductive Relation.
  \end{definition}
\begin{definition}[Computable Function]
  \label{dfn:MOList.Computation.next}
  \lean{MOList.Computation.next}
  \leanok
  Function computing the next state of the system.
  \end{definition}

\paragraph*{Sorted Properties}

%%%%%%%%%%%%%%%%%%%%%%%%%%%%%%%%%%%%%%%%
\begin{theorem}
  \label{thm:MOList.Properties.key_sorted}
  \lean{MOList.Properties.key_sorted}
  \leanok
  \uses{dfn:MUList.State.ListToken, dfn:ComputationalRelation.StateMachine.FromInit}
  \end{theorem}
\begin{proof}
  \proves{thm:MOList.Properties.key_sorted}
  \uses{thm:Tests.unordered_list_tests}
  We employed the Plausible Framework to test the above property.
  \end{proof}
%%%%%%%%%%%%%%%%%%%%%%%%%%%%%%%%%%%%%%%%

%%%%%%%%%%%%%%%%%%%%%%%%%%%%%%%%%%%%%%%%
\begin{theorem}
  \label{thm:MOList.Properties.all_diff_keys}
  \lean{MOList.Properties.all_diff_keys}
  \leanok
  \uses{dfn:MUList.State.ListToken, dfn:ComputationalRelation.StateMachine.FromInit}
  Keys are unique.
  \end{theorem}

\begin{proof}
  \proves{thm:MOList.Properties.all_diff_keys}
  \uses{thm:Tests.ordered_list_tests}
  We can deduce this from keys being sorted.
  %
  However, we employed the Plausible Framework to test the above property.
  \end{proof}
%%%%%%%%%%%%%%%%%%%%%%%%%%%%%%%%%%%%%%%%

%%%%%%%%%%%%%%%%%%%%%%%%%%%%%%%%%%%%%%%%
\begin{theorem}
  \label{thm:MOList.Properties.unique_tokens}
  \lean{MOList.Properties.unique_tokens}
  \leanok
  \uses{dfn:MUList.State.ListToken, dfn:ComputationalRelation.StateMachine.FromInit}
  There is (at most) only one token for each key.
  \end{theorem}
%%%%%%%%%%%%%%%%%%%%%%%%%%%%%%%%%%%%%%%%

%%%%%%%%%%%%%%%%%%%%%%%%%%%%%%%%%%%%%%%%
\begin{theorem}
  \label{thm:MOList.Properties.not_anchor_self}
  \lean{MOList.Properties.not_anchor_self}
  \leanok
  \uses{dfn:MUList.State.ListToken, dfn:ComputationalRelation.StateMachine.FromInit}
  No token is anchor of itself.
  \end{theorem}

\begin{proof}
  \proves{thm:MOList.Properties.not_anchor_self}
  \uses{thm:Tests.ordered_list_tests}
  Tested with plausible.
  \end{proof}

%%%%%%%%%%%%%%%%%%%%%%%%%%%%%%%%%%%%%%%%
\begin{theorem}
  \label{thm:MOList.Properties.neq_anchor_ids}
  \lean{MOList.Properties.neq_anchor_ids}
  \leanok
  \uses{dfn:MUList.State.ListToken, dfn:ComputationalRelation.StateMachine.FromInit}
  Given a token and its anchor, then the two are different.
  \end{theorem}
\begin{proof}
  \proves{thm:MOList.Properties.neq_anchor_ids}
  \uses{thm:MOList.Properties.not_anchor_self}
  Tokens are not anchors of themselves.
  \end{proof}
%%%%%%%%%%%%%%%%%%%%%%%%%%%%%%%%%%%%%%%%

\paragraph*{Membership}
%

%%%%%%%%%%%%%%%%%%%%%%%%%%%%%%%%%%%%%%%%
\begin{theorem}
  \label{thm:MOList.Properties.non_membership_root_and_first}
  \lean{MOList.Properties.non_membership_root_and_first}
  \leanok
  \uses{dfn:MUList.State.ListToken, dfn:ComputationalRelation.StateMachine.FromInit}
  All keys less than the key of the first element are not in the list.
  \end{theorem}
\begin{proof}
  \proves{thm:MOList.Properties.non_membership_root_and_first}
  \uses{thm:MOList.Properties.key_sorted}
  Consequences of keys being sorted.
  \end{proof}

\begin{theorem}
  \label{thm:MOList.Properties.non_membership_after_last}
  \lean{MOList.Properties.non_membership_after_last}
  \leanok
  \uses{dfn:MUList.State.ListToken, dfn:ComputationalRelation.StateMachine.FromInit}
  Keys greater than the last element are not in the list.
  \end{theorem}
\begin{proof}
  \proves{thm:MOList.Properties.non_membership_after_last}
  \uses{thm:MOList.Properties.key_sorted}
  Consequences of keys being sorted.
  \end{proof}

\begin{theorem}
  \label{thm:MOList.Properties.non_membership_witness}
  \lean{MOList.Properties.non_membership_witness}
  \leanok
  \uses{dfn:MUList.State.ListToken, dfn:ComputationalRelation.StateMachine.FromInit}
  Given a token whose key is \(k\) and link \(nk\),
  all keys \(q \in (k,nk)\) are not members of the list.
  \end{theorem}

\begin{proof}
  \proves{thm:MOList.Properties.non_membership_witness}
  \uses{thm:MOList.Properties.key_sorted}
  Consequences of keys being sorted.
  \end{proof}
%%%%%%%%%%%%%%%%%%%%%%%%%%%%%%%%%%%%%%%%

%%%%%%%%%%%%%%%%%%%%%%%%%%%%%%%%%%%%%%%%
\begin{theorem}
  \label{thm:MOList.Properties.non_membership_ne_member}
  \lean{MOList.Properties.non_membership_ne_member}
  \leanok
  \uses{dfn:MUList.State.ListToken, dfn:ComputationalRelation.StateMachine.FromInit}
  Midgard Ordered lists \emph{not\_membership} operation is correct.
  \end{theorem}

\begin{proof}
  \proves{thm:MOList.Properties.non_membership_ne_member}
  \uses{thm:MOList.Properties.non_membership_witness,thm:MOList.Properties.non_membership_after_last,thm:MOList.Properties.non_membership_root_and_first}
  Case analysis on the witness proving non\_membership.
  %
  We have three cases: witness is the first, the last, or an element between the
  previous two.
  %
  On each case, we can apply on of the previous theorems.
  \end{proof}
%%%%%%%%%%%%%%%%%%%%%%%%%%%%%%%%%%%%%%%%

%%%%%%%%%%%%%%%%%%%%%%%%%%%%%%%%%%%%%%%%
\begin{theorem}
  \label{thm:MOList.Properties.is_member_not_non_member}
  \lean{MOList.Properties.is_member_not_non_member}
  \leanok
  \uses{dfn:MUList.State.ListToken, dfn:ComputationalRelation.StateMachine.FromInit}
  If we have proof of membership, i.e \(\mathsf{is\_member(key,wit)}\), then
  it is not possible to have prove of
non-membership:\(\not \mathsf{is\_non\_member(key,wit)}\).
  \end{theorem}

\begin{theorem}
  \label{thm:MOList.Properties.non_member_is_not_member}
  \lean{MOList.Properties.non_member_is_not_member}
  \leanok
  \uses{dfn:MUList.State.ListToken, dfn:ComputationalRelation.StateMachine.FromInit}
  If we have proof of non-membership, i.e \( \mathsf{is\_non\_member(key,wit)}\)
then it is not possible to have prove of membership:
\(\not \mathsf{is\_member(key,wit)}\).
  \end{theorem}

Both properties above were tested using Plausible see Properties~\ref{dfn:Tests.ordered_list_tests}.

%%%%%%%%%%%%%%%%%%%%%%%%%%%%%%%%%%%%%%%%
The rest of the properties are usefull for proving other modules.

\paragraph*{Insertion}

\begin{theorem}
  \label{thm:MOList.Properties.insert_non_member}
  \lean{MOList.Properties.insert_non_member}
  \leanok
  \uses{dfn:MUList.State.ListToken, dfn:MOList.Computation.next}
  It is always possible to insert an element for which we have a witness stating
that there is not in the list.
  \end{theorem}

\begin{theorem}
  \label{thm:MOList.Properties.inserted_non_member}
  \lean{MOList.Properties.inserted_non_member}
  \leanok
  \uses{dfn:MUList.State.ListToken, dfn:MOList.Relation.Rel}
  If we just inserted an element, then such element was not in the list before
inserting.
  \end{theorem}

\begin{theorem}
  \label{thm:MOList.Properties.not_insert_member}
  \lean{MOList.Properties.not_insert_member}
  \leanok
  \uses{dfn:MUList.State.ListToken, dfn:MOList.Relation.Rel}
  It is not possible to insert an element that it already in the list.
  \end{theorem}

\begin{theorem}
  \label{thm:MOList.Properties.insert_diff_new}
  \lean{MOList.Properties.insert_diff_new}
  \leanok
  \uses{dfn:MUList.State.ListToken, dfn:MOList.Relation.Rel}
  After inserting an element, the element is in the list.
  \end{theorem}

\begin{theorem}
  \label{thm:MOList.Properties.insert_new_iff}
  \lean{MOList.Properties.insert_new_iff}
  \leanok
  \uses{dfn:MUList.State.ListToken, dfn:MOList.Relation.Rel}
  After inserting an element, an element belongs to the resulting list if it is
either the just inserted node or another belonging to the previous list.
  \end{theorem}

\begin{theorem}
  \label{thm:MOList.Properties.insert_diff_same_old}
  \lean{MOList.Properties.insert_diff_same_old}
  \leanok
  \uses{dfn:MUList.State.ListToken, dfn:MOList.Relation.Rel}
  After inserting, all the members of the previous list are members of the new list.
  \end{theorem}

\begin{theorem}
  \label{thm:MOList.Properties.insert_diff_same_old}
  \lean{MOList.Properties.insert_diff_same_old}
  \leanok
  \uses{dfn:MUList.State.ListToken, dfn:MOList.Relation.Rel}
  After inserting, all members of the next list, but the just inserted one,
belongs to the list before inserting.
  \end{theorem}

\paragraph*{Remove Operations.}

\begin{theorem}
  \label{thm:MOList.Properties.EnabledRemove'}
  \lean{MOList.Properties.EnabledRemove'}
  \leanok
  \uses{dfn:MUList.State.ListToken, dfn:MOList.Relation.Rel}
  We can always remove elements by providing the token and its anchors.
  \end{theorem}

\begin{theorem}
  \label{thm:MOList.Properties.OrdRemove_UordRemove}
  \lean{MOList.Properties.OrdRemove_UordRemove}
  \leanok
  \uses{dfn:MUList.State.ListToken, dfn:MOList.Relation.Rel}
  Ordered List action removes is implementing using Unorederd List action
remove.
  \end{theorem}

\begin{theorem}
  \label{thm:MOList.Properties.AfterRemove_TId_not_member}
  \lean{MOList.Properties.AfterRemove_TId_not_member}
  \leanok
  \uses{dfn:MUList.State.ListToken, dfn:MOList.Relation.Rel}
  After removing an element from the list, it is not possible to use the token
to prove membership.
  \end{theorem}
  
\begin{theorem}
  \label{thm:MOList.Properties.AfterRemoveKey_not_member}
  \lean{MOList.Properties.AfterRemoveKey_not_member}
  \leanok
  \uses{dfn:MUList.State.ListToken, dfn:MOList.Relation.Rel}
  After removing an element from the list, the key of the removed token is not in the list.
  \end{theorem}

\begin{theorem}
  \label{thm:MOList.Properties.OrdRemoves_members_back}
  \lean{MOList.Properties.OrdRemoves_members_back}
  \leanok
  \uses{dfn:MUList.State.ListToken, dfn:MOList.Relation.Rel}
  After removing an element, if an element is a member of the new list, then it
belonged also to the list before removing an element.
  \end{theorem}

\subsection{Plausible Generation}
\label{sec:list:plausible}

To quickly test properties about Midgard lists, we implemented Plausible
generators.
%
Plausible is the Lean4 property-based framework.
%
The generators simple keep up a simplified view of the internal state and
generate possible values on demand.
%
For a more detailed explanation, we refer the reader to the DocBuild
documentation.

\begin{theorem}
  \label{thm:Tests.view_compliance_tests }
  \lean{Tests.view_compliance_tests }
  \leanok
  \uses{dfn:MUList.State.ListToken, dfn:MUList.State.Actions}
  We tested that the views generated match the internal states.
  \end{theorem}

\begin{theorem}
  \label{thm:Tests.unordered_list_tests} 
  \lean{Tests.unordered_list_tests} 
  \leanok
  \uses{dfn:MUList.State.ListToken, dfn:MUList.State.Actions}
  Battery test for Midgard Unoredred Lists.
  \end{theorem}

%%%%%%%%%%%%%%%%%%%%%%%%%%%%%%%%%%%%%%%%

\begin{theorem}
  \label{thm:Tests.view_checks} 
  \lean{Tests.view_checks} 
  \leanok
  \uses{dfn:MUList.State.ListToken, dfn:MUList.State.Actions}
  Generated views are true to the internal states.
  \end{theorem}

\begin{theorem}
  \label{thm:Tests.ordered_list_tests}
  \lean{Tests.ordered_list_tests}
  \leanok
  \uses{dfn:MUList.State.ListToken, dfn:MUList.State.Actions}
  Battery test for MidgardOrdered Lists.
  \end{theorem}
% TODO Define {dfn:List.SortedCheckProp.sorted_key_list }
% TODO Define {dfn:Tests.ordered_list_tests }


%%FILE blueprint/src/modules/OperatorDirectory.tex
\section{Operator Directory}

Operator Directory is the Midgard module in charge of registering, activating
and retiring operators.
%
In this module, we operate three lists (or queues), handling operators three
possible states: 1) registered, 2) active, 3) retired.
%
Operators can register (and deregister) freely by placing (and taking) a
required bond.
%
The amount of the required bond is set as a Midgard parameter.
%
Registered operators can be activated after a pre-defined time period has passed.
%
Active operators can retire anytime and are the ones in charge of advancing the
state of Midgard, i.e. of pushing new L2-Midgard blocks.
%
Retired operators need to wait until their bonds are free to be claimed.
%
The bonds are free once all their pushed blocks are committed to the history of
Midgard and cannot be challenge.

Operators activation happens when a scheduler cycle ends, and the scheduler is
rewound.
%
Registered operators are eligible for activation after a waiting for a preset
waiting time.
%
The length of this waiting period is also a Midgard parameter.
%

Properties in this module come in two flavours: invariants and optimistic properties.
%
Invariants are classic properties of the system, see Section~\ref{sec:opdir:invariants}.
%
Optimistic Properties are not properties in the classical sense, optimistic properties are not invariants.

Optimistic properties when became properties when we reason about economic
incentives plus rational agents.
%
The system, the operator directory module in this case, accepts users to perform
actions leading the system to states denoted as fraudulent.
%
When these states are reached, then the system also accepts other users to
provide evidence of fraudulent actions and punishes the users leading the system
to such states.
%
For a graphical representation see the finite state machine representation, see
Figure~\ref{fig:opdir}.
%
We model these complex interactions in Section~\ref{sec:opdir:optimisticprops}.

\subsection{Plots}

The Operator directory module accepts fraudulent operations and also supports
the punishment of fraudulent operators.
%
In other words, it supports \emph{Optimistic Properties.}
%
Hence, we present here two finite state machine plots, one with all the possible
behaviours and another with the desired ones.

\begin{figure}
  \caption{Operator Directory Finite State Machine. Desired and Actual state space.}
  \includegraphics[width=\textwidth]{operatorsdirectory}
  \label{fig:opdir}
\end{figure}


\begin{figure}
  \caption{Operator Directory IO plot.}
  \includegraphics[width=\textwidth]{io.opdir}
  \label{fig:opdir}
\end{figure}

\subsection{Requires}

\paragraph{Unordered Simple Linked Lists:}
 %
 Registered operators queue is a FIFO.
 %
 Upon registration, Midgard takes note of the time of the Cardano transaction,
 through the upper transaction validity bound, and adds a fix waiting period.

\paragraph{Ordered Simple Linked Lists:}
\begin{itemize}
 \item Active operators queue. 
 \item Retired operators queue.
\end{itemize}

\subsection{Formalization}

\begin{definition}[Operator Directory]
  \label{dfn:OperatorDirectory.State.OperatorDirectory}
  \lean{OperatorDirectory.State.OperatorDirectory}
  We define operators directory as three lists: one unordered for registered
  operators, and two ordered for active and retired operators.
  \end{definition}

One way of defining the system could be through complex composition of the list
machines.
%
We opted for collapsing operations happening at the same time instead of using
products and sub-interfacing of machines.

\begin{definition}[State Transitions]
  \label{dfn:OperatorDirectory.StateMachine.State.OperActions}
  \lean{OperatorDirectory.State.OperActions}\leanok
  Operators can perform the following actions:
  \begin{description}
      \item[Activate] activates registered operators 
      \item[Retire] retires active operators 
      \item[Register] registers an operator by placing a bond, the operator is the one forging the transaction.
	  \item[Deregister] deregisters an operator (if possible)
      \item[RemoveDuplicateSlashBond] removes fraudulent operator
	  \item[RecoverBond] returns (if possible) the bond to a retired operator
	  \item[ActRemoveOperatorBad] removes a fraudulent active operators,
		  used by BadSettlement and BadBlock
	  \item[UpdateBondHold] updates bond holding time, used by New State
		  and New Settlement
	  \item[RetRemove] removes fraudulent retired operators, used by Bad
		  State and Settlement.
     \end{description}
  \end{definition}

\begin{definition}[Operators Directory Inductive Relation]
  %% Def and Lean
  \label{dfn:OperatorDirectory.Relation}
  \lean{OperatorDirectory.Relation}\leanok
  %% Uses
  \uses{dfn:MidgardParameters,dfn:InputInfo,dfn:OperatorDirectory.StateMachine.State.OperActions,dfn:OperatorDirectory.State.OperatorDirectory}
  %%
  Defining transitions based on Midgards parameters, input information from the
	Cardano L1 transaction, and information provided by the user.
  \end{definition}

\begin{theorem}[Operators Dir Relation is Computational]
  %%
  \label{thm:OperatorDirectory.Computational.ComputationalRelation}
  \leanok
  \lean{OperatorDirectory.Computational.ComputationalRelation}
  \uses{dfn:computational_relation,dfn:OperatorDirectory.Relation}
  %%
  There is a function realizing the computation of the Operators Directory
  inductive relation.
  \end{theorem}

\subsection{Properties}

\subsubsection{Invariants}\label{sec:opdir:invariants}

We list here the more interesting properties first, and then a list of technical
properties.

\paragraph{Registered Operators sorted incr by data (registered time).}
%
Upon registration, registration operator's time is defined by the transaction
registering them.
%
Time is increasing in Cardano and new operators are prepended to the list, so
their registration time is decreasing.

\redflag{Time in Cardano increases, but the transaction validity bounds are
chosen by users within certain restrictions. However, this brakes the assumption
that time (observed by Midgard) increases, rendering this property useless.}

Midgard's specification assumes that the last node in the registered operators
queue is the oldest node in the list.
%
It is the oldest node, but they are not ordered by time, so there may be
eligible nodes for activation even if the last one is not.

\redflag{The Midgard specification assumes that this property holds. For
example, when checking if there is a registered operator can be activated.}
\martin{We raised this problem, and the people from Anastesia labs said that it
was easy to fix. We only need to have a clear criteria for activating
registered operators.}

\paragraph{Register Operators}

Anyone can register, assuming the operator can generate a valid transaction.

\begin{theorem}[Anyone can register]
  \uses{dfn:MidgardParameters,dfn:OperatorDirectory.State.OperatorDirectory}
  \lean{OperatorDirectory.Properties.Registered.AnyoneCanRegister}
  \label{thm:OperatorDirectory.Properties.Registered.AnyoneCanRegister}
  \leanok
  Any operator can register themselves through a transaction.
  \end{theorem}

\begin{proof}
  \proves{thm:OperatorDirectory.Properties.Registered.AnyoneCanRegister}
  Proof is very simple, we show that operator registering action is always
enabled providing a valid Cardano transaction.
  \end{proof}

\begin{theorem}[Activating Registered Operators]
  \uses{dfn:MidgardParameters,dfn:OperatorDirectory.State.OperatorDirectory}
  \lean{OperatorDirectory.Properties.Registered.RegisteredBecomeActive}
  \label{thm:OperatorDirectory.Properties.Registered.RegisteredBecomeActive}
  \leanok
Registered nodes can be activated if their activation time has passed.
  \end{theorem}
\begin{theorem}
  \uses{dfn:MidgardParameters,dfn:OperatorDirectory.State.OperatorDirectory}
  \lean{OperatorDirectory.Properties.Registered.RegisterDeregister}
  \label{thm:OperatorDirectory.Properties.Registered.RegisterDeregister}
  \leanok
Registered nodes can deregisterd \textbf{themselves}.
  \end{theorem}

\paragraph{Active Operators}

\begin{theorem}
  \uses{dfn:MidgardParameters,dfn:OperatorDirectory.State.OperatorDirectory}
  \lean{OperatorDirectory.Properties.ActiveOpDir.ActiveNodesRetire}
  \label{thm:OperatorDirectory.Properties.ActiveOpDir.ActiveNodesRetire}
  \leanok
Active nodes can be retired by \textbf{anyone}.
  \end{theorem}
\martin{This is a but, we reported it.}

\paragraph*{Active and Retired}

\begin{theorem}
  \uses{dfn:MidgardParameters,dfn:OperatorDirectory.State.OperatorDirectory}
  \lean{OperatorDirectory.Properties.NonRetiredActive}
  \label{thm:OperatorDirectory.Properties.NonRetiredActive}
  \leanok
  Active and retired lists are disjoint.
  \end{theorem}

%%%%%%%%%%%%%%%%%%%%%%%%%%%%%%%%%%%%%%%%%%%%%%%%%%%%%%%%%%%%%%%%%%%%%%%%%%%%%%%%


\paragraph{Sub-Tracing}
%
Every state reachable by the Operator Directory machine has three lists that are
reachable states by the corresponding list machine.

\begin{theorem}[Registered Reachable List]
  \leanok
  \lean{OperatorDirectory.Properties.ListConsistency.registered_reachable}
  \label{thm:OperatorDirectory.Properties.ListConsistency.registered_reachable}
  %%
  The registered list is a valid (reachable) state.
\end{theorem}
\begin{theorem}[Active Reachable List]
  \leanok
  \lean{OperatorDirectory.Properties.ListConsistency.active_reachable}
  \label{thm:OperatorDirectory.Properties.ListConsistency.active_reachable}
  %%
  The registered list is a valid (reachable) state.
\end{theorem}
\begin{theorem}[Retired Reachable List]
  \leanok
  \lean{OperatorDirectory.Properties.ListConsistency.retired_reachable}
  \label{thm:OperatorDirectory.Properties.ListConsistency.retired_reachable}
  %%
  The registered list is a valid (reachable) state.
\end{theorem}

\paragraph{Actions}

When composing operators directory state machine with others, we need to
explore the different properties we get from triggering actions.
%
Here we list some of those properties.


\paragraph*{Activating Properties.}

\begin{theorem}[Activate Postcondition]
  \leanok
  \lean{OperatorDirectory.Properties.Actions.activate_props}
  \uses{dfn:OperatorDirectory.Relation}
  \label{thm:OperatorDirectory.Properties.Actions.activate_props}

  After activating a node: registered node is removed from the registered
		queue, and the operator now belongs to the activated queue sorted
		(decr) by operator id.
\end{theorem}


\begin{theorem}
  \leanok
  \lean{OperatorDirectory.Properties.Actions.activate_insert_activates}
  \label{thm:OperatorDirectory.Properties.Actions.activate_insert_activates}
  When activating an operator, the operator is inserted in the active operator list.
\end{theorem}

\begin{theorem}
  \leanok
  \lean{OperatorDirectory.Properties.Actions.activate_register_removes}
  \label{thm:OperatorDirectory.Properties.Actions.activate_register_removes}
  When activating an operator, the operator is removed from the registered operator list.
\end{theorem}
\begin{theorem}
  \leanok
  \lean{OperatorDirectory.Properties.Actions.activate_register_not_self}
  \label{thm:OperatorDirectory.Properties.Actions.activate_register_not_self}
  Operators cannot activate using themselves as anchor.
\end{theorem}
\begin{theorem}
  \leanok
  \lean{OperatorDirectory.Properties.Actions.activate_register_non_retired}
  \label{thm:OperatorDirectory.Properties.Actions.activate_register_non_retired}
  If we can activate an operator then we know the operator is not a member of
the retired list.
\end{theorem}
\begin{theorem}
  \leanok
  \lean{OperatorDirectory.Properties.Actions.activate_same_retired}
  \label{thm:OperatorDirectory.Properties.Actions.activate_same_retired}
  Activating an operator does not modify the retired operator list.
\end{theorem}


%%%
\paragraph*{Retiring Properties.}

\begin{theorem}
  \leanok
  \lean{OperatorDirectory.Properties.Actions.retired_active_substep}
  \label{thm:OperatorDirectory.Properties.Actions.retired_active_substep}
  %
  Retiring an operator removes such operator from the activated operator list.
  \end{theorem}

\begin{theorem}
  \leanok
  \lean{OperatorDirectory.Properties.Actions.retired_retired_step}
  \label{thm:OperatorDirectory.Properties.Actions.retired_retired_step}
  %
  Retiring an operator moves the operator to the retired list maintaining the
same information.
  \end{theorem}

\begin{theorem}[Retiring keeps registered]
  \leanok
  \lean{OperatorDirectory.Properties.Actions.retired_registered}
  \uses{dfn:OperatorDirectory.Relation}
  \label{thm:OperatorDirectory.Properties.Actions.retired_registered}
  Retiring operators do not modify registered operators queue.
  \end{theorem}

\paragraph*{Register Properties.}

\begin{theorem}[Register keeps active]
  \leanok
  \lean{OperatorDirectory.Properties.Actions.register_same_active}
  \uses{dfn:OperatorDirectory.Relation}
  \label{thm:OperatorDirectory.Properties.Actions.register_same_active}
  Registering operators do not modify the active operator list.
  \end{theorem}
\begin{theorem}[Register keeps retired]
  \leanok
  \lean{OperatorDirectory.Properties.Actions.register_same_retired}
  \uses{dfn:OperatorDirectory.Relation}
  \label{thm:OperatorDirectory.Properties.Actions.register_same_retired}
  Registering operators do not modify the retired operator list.
  \end{theorem}

\begin{theorem}[Post condition after registering operators]
  \leanok
  \lean{OperatorDirectory.Properties.Actions.register_register}
  \uses{dfn:OperatorDirectory.Relation}
  \label{thm:OperatorDirectory.Properties.Actions.register_register}
  After operator registration, we have that all registered nodes are kept the same (but root?.)
  %
  Fresh minted token did not exist before registration.
  %
  Fresh minted token proves membership of the registered operator.
  \end{theorem}

\begin{theorem}[Operators are Prepended.]
  \leanok
  \lean{OperatorDirectory.Properties.Actions.register_prepend}
  \label{thm:OperatorDirectory.Properties.Actions.register_prepend}
  \uses{dfn:OperatorDirectory.Relation}
  Operators are prepended (added to the end) of the registered operator list.
  \end{theorem}


\paragraph*{Deregister Properties.}

\begin{theorem}[Deregister keeps active]
  \leanok
  \lean{OperatorDirectory.Properties.Actions.deregister_same_active}
  \label{thm:OperatorDirectory.Properties.Actions.deregister_same_active}
  \uses{dfn:OperatorDirectory.Relation}
  Deregistering operators do not modify the active operator list.
  \end{theorem}
\begin{theorem}[Deregister keeps retired]
  \leanok
  \lean{OperatorDirectory.Properties.Actions.deregister_same_retired}
  \label{thm:OperatorDirectory.Properties.Actions.deregister_same_retired}
  \uses{dfn:OperatorDirectory.Relation}
  Deregistering operators do not modify the retired operator list.
  \end{theorem}

\begin{theorem}[Post condition after deregistering operators.]
  \leanok
  \lean{OperatorDirectory.Properties.Actions.deregister_registered}
  \uses{dfn:OperatorDirectory.Relation}
  \label{thm:deregister_register}
  Deregistering an operator removes the registered node, updates the anchor node
properly and keeps the rest of the nodes untouched.
  \end{theorem}


\paragraph*{Remove Duplicate Properties.}

\begin{theorem}[Remove duplicates keeps active]
  \leanok
  \lean{OperatorDirectory.Properties.Actions.removeduplicatereg_same_active}
  \label{thm:OperatorDirectory.Properties.Actions.removeduplicatereg_same_active}
  \uses{dfn:OperatorDirectory.Relation}
  Removing duplicated operators do not modify the active operator list.
  \end{theorem}
\begin{theorem}[Remove duplictes keeps retired]
  \leanok
  \lean{OperatorDirectory.Properties.Actions.removeduplicatereg_retired}
  \label{thm:OperatorDirectory.Properties.Actions.removeduplicatereg_retired}
  \uses{dfn:OperatorDirectory.Relation}
  Removing duplicated operators do not modify the retired operator list.
  \end{theorem}

\begin{theorem}[Remove duplicated removes one registered operator]
  \leanok
  \lean{OperatorDirectory.Properties.Actions.removeduplicatereg_removes_dup_registered}
  \label{thm:OperatorDirectory.Properties.Actions.removeduplicatereg_removes_dup_registered}
  \uses{dfn:OperatorDirectory.Relation}

  Removing duplicated registered operators removes a node from the registered
operators queue.

  \end{theorem}

\paragraph*{Active Remove Properties.}
\begin{theorem}
  \leanok
  \lean{OperatorDirectory.Properties.Actions.active_remove_bad_registered}
  \label{thm:OperatorDirectory.Properties.Actions.active_remove_bad_registered}
  \uses{dfn:OperatorDirectory.Relation}
  The active operator list is not modified.
  \end{theorem}
\begin{theorem}
  \leanok
  \lean{OperatorDirectory.Properties.Actions.retired_remove_bad_registered}
  \label{thm:OperatorDirectory.Properties.Actions.retired_remove_bad_registered}
  \uses{dfn:OperatorDirectory.Relation}
  The retired operator list is not modified.
  \end{theorem}

\begin{theorem}[Active Remove Bad Active]
  \leanok
  \lean{OperatorDirectory.Properties.Actions.active_remove_bad_active_keys}
  \label{thm:OperatorDirectory.Properties.Actions.active_remove_bad_active_keys}
  \uses{dfn:OperatorDirectory.Relation}
  Active remove bad active operator removes such operator.
  \end{theorem}

\paragraph*{Update Bonds properties.}

\begin{theorem}[Post condition after Update Bonds action.]
  \leanok
  \lean{OperatorDirectory.Properties.Actions.update_bond_hold_active_keys}
  \label{thm:OperatorDirectory.Properties.Actions.update_bond_hold_active_keys}
  \uses{dfn:OperatorDirectory.Relation}
  After successfully updating the bonds of an operator, we get:
  all other active members remain the same, the transaction internal validity
upper bound is defined, the active operator exists, keeps existing after the
update, and the state is properly updated.
  \end{theorem}

\begin{theorem}
  \leanok
  \lean{OperatorDirectory.Properties.Actions.update_bond_hold_same_act_members}
  \label{thm:OperatorDirectory.Properties.Actions.update_bond_hold_same_act_members}
  \uses{dfn:OperatorDirectory.Relation}
  Updating bonds do not change the membership status in the active operator list.
  \end{theorem}

\begin{theorem}
  \leanok
  \lean{OperatorDirectory.Properties.Actions.update_bond_hold_same_registered}
  \label{thm:OperatorDirectory.Properties.Actions.update_bond_hold_same_registered}
  \uses{dfn:OperatorDirectory.Relation}
  Updating bonds do not change the registered operator list.
  \end{theorem}
\begin{theorem}
  \leanok
  \lean{OperatorDirectory.Properties.Actions.update_bond_hold_same_retired}
  \label{thm:OperatorDirectory.Properties.Actions.update_bond_hold_same_retired}
  \uses{dfn:OperatorDirectory.Relation}
  Updating bonds do not change the retired operator list.
  \end{theorem}

\paragraph*{Recover Bond Properties.}

\begin{theorem}
  \leanok
  \lean{OperatorDirectory.Properties.Actions.recover_bond_same_active}
  \label{thm:OperatorDirectory.Properties.Actions.recover_bond_same_active}
  \uses{dfn:OperatorDirectory.Relation}
  Recovering bonds do not change the active operator list.
  \end{theorem}
\begin{theorem}
  \leanok
  \lean{OperatorDirectory.Properties.Actions.recover_bond_same_registered}
  \label{thm:OperatorDirectory.Properties.Actions.recover_bond_same_registered}
  \uses{dfn:OperatorDirectory.Relation}
  Recovering bonds do not change the registered operator list.
  \end{theorem}
\begin{theorem}
  \leanok
  \lean{OperatorDirectory.Properties.Actions.recover_bond_retired}
  \label{thm:OperatorDirectory.Properties.Actions.recover_bond_retired}
  \uses{dfn:OperatorDirectory.Relation}
  Recovering bonds removes the corresponding node from the retired operator list.
  \end{theorem}

\paragraph*{Remove Bad Retired Operator}
\begin{theorem}
  \leanok
  \lean{OperatorDirectory.Properties.Actions.retired_remove_same_active}
  \label{thm:OperatorDirectory.Properties.Actions.retired_remove_same_active}
  \uses{dfn:OperatorDirectory.Relation}
  The active operator list does not change.
  \end{theorem}
\begin{theorem}
  \leanok
  \lean{OperatorDirectory.Properties.Actions.retired_remove_same_registered}
  \label{thm:OperatorDirectory.Properties.Actions.retired_remove_same_registered}
  \uses{dfn:OperatorDirectory.Relation}
  The registered operator list does not change.
  \end{theorem}
\begin{theorem}
  \leanok
  \lean{OperatorDirectory.Properties.Actions.retired_remove_retired}
  \label{thm:OperatorDirectory.Properties.Actions.retired_remove_retired}
  \uses{dfn:OperatorDirectory.Relation}
  The corresponding operator is removed from the retired operator list.
  \end{theorem}


\subsubsection{Optimistic Properties}\label{sec:opdir:optimisticprops}
%
Optimistic Properties are different than normal invariants.
%
These properties are not properties in the classical sense,
but become properties when reasoning using economic incentives and rational agents.
%
The system accepts actions from agents leading to invalid states, but when
detected, these agents are severily penalized.
%
The system does not perform any check avoiding unwanted states, but needs to
guarantee that \textbf{only and always} fraudulent agents are penalized.
%
Moreover, the system has to be able to recover from fraudulent states and
continue operating from a non-fraudulent state.
%
The economic reasoning follows from the existence of severe penalizations of
fraudulent agents and rewards for honest agents watching the evolution of the
system.
%
Economic analyses are left outside this formalization.

Onchain code benefits from these Optimistic Properties.
%
There are less operations performed onchain, and in the rare case that someone
does something wrong, we trigger checks and operations penalizing the
fraudulent agents.
%
However, \emph{rational agents} will not take actions leading to penalizations.
%
The mere existence of detecting and penalizing agents is a deterrent system.
%
Our job here, as members of the formal methods team, is to show that such agents
will always be caught and penalized.

There are three optimistic properties in the Operator Directory module:
%
\begin{enumerate}
  \item No Duplicated Registered Nodes
  \item No operator is Registered and Active at the same time
  \item No operator is Registered and Retired at the same time
  \end{enumerate}

We know from before~(see
Theorem~\ref{thm:OperatorDirectory.Properties.NonRetiredActive}) that there are
no active and retired operators.

To prove optimistic invariants, we prove that:
%
\begin{enumerate}
    \item Agents can \emph{detect} invalid states.
    \item Agents can \emph{generate actions} to \emph{correct} it.
    \item Assuming agents can \emph{submit} such actions to the Midgard, the
system penalizes fraudulent operators.
  \end{enumerate}
  %
We assume agents can always submit operations to the Cardano mainnet.

When put altogether, the reasoning behind this that: operators can always detect
fraudulent states, they can collect the necessary evidence the system requires,
and submit it to the system so fraudulent operators are penalized.
%

We reason using concepts from \emph{alternating temporal logic}, a logic
equipped with agentic notions: Midgard operators can collude to enforce certain
properties.


\begin{definition}[Register Queue Violation Detection]
  %%
  \label{dfn:OpsDir.OptimisticDetection.find_witness}
  \uses{dfn:OperatorDirectory.State.OperatorDirectory}
  \lean{OpsDir.OptimisticDetection.find_witness}
  \leanok
  %%
  Finds if there are two elements in the registered queue with the same key and
different times or there is an operator registered and active or registered and
retired.
%
  Here we assume agents have a valid state of the onchain data structure.
  \end{definition}

\begin{theorem}
  \uses{dfn:OpsDir.OptimisticDetection.find_witness,dfn:ComputationalRelation.StateMachine.FromInit,dfn:OperatorDirectory.Relation}
  \lean{OpsDir.OptimisticDetection.find_witness_valid}
  \label{thm:OpsDir.OptimisticDetection.find_witness_valid}
  \leanok
  The witness found generates a \emph{valid witness}.
  %
  It finds two different tokens with matching keys.
  \end{theorem}

\begin{theorem}
  \uses{dfn:OpsDir.OptimisticDetection.find_witness,dfn:ComputationalRelation.StateMachine.FromInit,dfn:OperatorDirectory.Relation}
  \lean{OpsDir.OptimisticDetection.find_witness_key}
  \label{thm:OpsDir.OptimisticDetection.find_witness_key}
  \leanok
  The witness found identified a registered operator to penalized.
  %
  In other words, the operator to penalize is a member of the registered operator list.
  \end{theorem}


The Midgard protocol, when an operator is fraudulent, the operator is removed
from the registered operator list.
%
To remove operators, we need a valid anchor.

\begin{theorem}
  \uses{dfn:OpsDir.OptimisticDetection.find_witness,dfn:ComputationalRelation.StateMachine.FromInit,dfn:OperatorDirectory.Relation}
  \lean{OpsDir.OptimisticDetection.find_witness_registered}
  \label{thm:OpsDir.OptimisticDetection.find_witness_registered}
  \leanok
  The witness found has a valid anchor for the operator the
system is penalizing.
  \end{theorem}

\begin{theorem}[Duplicate Registered Detection]
  \uses{dfn:OpsDir.OptimisticDetection.find_witness,dfn:ComputationalRelation.StateMachine.FromInit,dfn:OperatorDirectory.Relation}
  \lean{OpsDir.OptimisticDetection.find_witness_dup_reg}
  \label{thm:OpsDir.OptimisticDetection.find_witness_dup_reg}
  \leanok
  If the witness found is marked as registered duplication, then the witness has
two distinct tokens showing duplication of a registered operator and another
registered operator.
  \end{theorem}
\begin{theorem}[Duplicate Registered and Active Detection]
  \uses{dfn:OpsDir.OptimisticDetection.find_witness,dfn:ComputationalRelation.StateMachine.FromInit,dfn:OperatorDirectory.Relation}
  \lean{OpsDir.OptimisticDetection.find_witness_dup_act}
  \label{thm:OpsDir.OptimisticDetection.find_witness_dup_act}
  \leanok
  If the witness found is marked as registered and active duplication, then the witness has
two distinct tokens showing duplication of a registered operator and an active operator.
  \end{theorem}
\begin{theorem}[Duplicate Registered and Retired Detection]
  \uses{dfn:OpsDir.OptimisticDetection.find_witness,dfn:ComputationalRelation.StateMachine.FromInit,dfn:OperatorDirectory.Relation}
  \lean{OpsDir.OptimisticDetection.find_witness_dup_ret}
  \label{thm:OpsDir.OptimisticDetection.find_witness_dup_ret}
  \leanok
  If the witness found is marked as registered duplication, then the witness has
  two distinct tokens showing duplication of a registered operator and
  a retired operator.
  \end{theorem}

  
\begin{theorem}[No Witness Detected]
  \uses{dfn:OpsDir.OptimisticDetection.find_witness,dfn:ComputationalRelation.StateMachine.FromInit,dfn:OperatorDirectory.Relation}
  \lean{OpsDir.OptimisticDetection.not_finding_witness}
  \label{thm:OpsDir.OptimisticDetection.not_finding_witness}
  \leanok
  If no witnesses are detected, then there are no duplication of registered operators.
  \end{theorem}

The above theorems guarantee us that agents monitoring the onchain state can
detect and generate a valid witness to submit to the Midgard protocol.
%
Next step is to generate a valid action at the Midgard level.

\begin{definition}[Proof Generation]
  \label{dfn:OpsDir.OptimisticProof.GenerateFraudProof}
  \lean{OpsDir.OptimisticProof.GenerateFraudProof}
  \leanok
  \uses{dfn:OperatorDirectory.Relation}
  Takes the witness generated by the condition checker and forms a valid Midgard Operator Directory action.
  \end{definition}

We show that when the condition checker generates a witness,
anyone can submit a fraudproof.
%
Anyone that can generate an acceptable trasaction context, i.e. submit an
L1-transaction.

\begin{theorem}[Enabled Register Fraudulent Witness Submission]
  %% MData
  \label{thm:OpsDir.OptimisticProof.EnabledFPSubmission}
  \lean{OpsDir.OptimisticProof.EnabledFPSubmission}
  \leanok
  \uses{dfn:OpsDir.OptimisticDetection.find_witness,dfn:MidgardParameters,dfn:InputInfo,dfn:OperatorDirectory.Relation}
  %%
  If we find a witness(see Definition~\ref{dfn:OpsDir.OptimisticDetection.find_witness}) and we can submit a
transaction to L1, then we can generate a \emph{valid} action
(using Definition~\ref{dfn:OpsDir.OptimisticProof.GenerateFraudProof}) having
Midgard penalizing the duplicated registered operator.
  \end{theorem}

We prove that if we execute the corrective measure, we get into a
different state which cannot be invalid by the same witness.
%

\begin{theorem}[Not same Witness]
  \label{thm:OpsDir.OptimisticProof.AfterFP}
  \uses{dfn:OpsDir.OptimisticDetection.find_witness, dfn:OperatorDirectory.Relation, dfn:OpsDir.OptimisticProof.GenerateFraudProof}
  \lean{OpsDir.OptimisticProof.AfterFP}
  \leanok
  %%
  Let \(s_{i},s_{o}\) be two operators directory such that
  \(s_{o}\) is the resulting state after slashing a fraudulent operator detected in \(s_{i}\).
  %
  Then, we won't find the same witness of fraudulent
condition~(Def.~\ref{dfn:OpsDir.OptimisticDetection.find_witness}).
  \end{theorem}

Finally, we have that if the condition returns no witness, we cannot submit
false fraudproof.
%
We express this property as a \emph{disabled} action.

\begin{theorem}[Disabled Duplicate Submission]
  \lean{OpsDir.OptimisticProof.DisableFPSubmittion}
  \label{thm:OpsDir.OptimisticProof.DisableFPSubmittion}
  \leanok
  \uses{dfn:MidgardParameters,dfn:OperatorDirectory.State.OperatorDirectory,dfn:InputInfo,dfn:OpsDir.OptimisticDetection.find_witness,dfn:OperatorDirectory.Relation}

  If when we check, we do not find any duplicates in the registered operators
queue, nor intersections with the retired and active queues, then  in the
current state, Midgard does not accept the removal of registered operators
accused of being duplicated.
  \end{theorem}


%%FILE blueprint/src/modules/Parameters.tex
\section{Midgard Parameters}

Midgard parameters module.

\begin{definition}{Parameters}
\label{dfn:MidgardParameters}
\lean{MidgardParameters}
\leanok
Record defining each parameter of a given Midgard instance.
\end{definition}


%%FILE blueprint/src/modules/ProofProtocol.tex
\section{Proof Protocol}

There are several modules interacting to form the \emph{Proof Protocol.}
%
First, the fraud proof catalog.
%
Second, we have the \emph{Computation Thread}, splitting potentially big
computations into simpler ones.
%
Third, we have the \emph{fraud proof tokens}, representing the correct
formulation of a fraud proof.

Once we have a fraud proof token, we can send it to the corresponding modules
and penalize fraudulent agents.

The \emph{Computation Thread} is implemented, but not used, see
Section~\ref{sec:computation_threads}.
%
We focused on checking the logic of fraud proofs and codified them directly in
Lean.

\subsection{Fraud Proof Tokens}
%
This module is just a collection of tokens.
%
The mere existence of this token is proof enough of fraud when presented at
other modules.
%
The correct generation is critical to the correct operation of Midgard.
%
Anyone should be able to generate them following the computation thread
mechanism.

\begin{definition}[FraudProof Tokens State]
  \lean{FraudProofTokens.State.FraudProofTokens}
  \leanok
  \label{dfn:FraudProofTokens.State.FraudProofTokens}
  The module state is just a collection of tokens indicating the hash of a
fraudulent block. 
  \end{definition}

\begin{definition}[Relation]
  \lean{FraudProofTokens.Computation.Relation}
  \leanok
  \label{dfn:FraudProofTokens.State.FraudProofTokens}
  There is only one action minting fraudproof tokens.
  %
  The only condition is to provide a successfully finished computation thread.
\end{definition}

We opted to define Lean4 interpretations of each fraud-proof described in the
Midgard specification.
%
\begin{definition}[Fraud Proofs]
  \lean{ProofProtocol.Catalogue.FraudProof}
  \label{dfn:ProofProtocol.Catalogue.FraudProof}
  \leanok
  Inductive type describing each fraudproof described in the Midgard specification.
\end{definition}

By using the above data-type, we skip the execution of computation threads.
%
Not implementing computation threads, means that we may be executing
computations that are too big to execute in one transaction in Cardano.

\subsection{Computation Threads}\label{sec:computation_threads}

Computation Threads are implemented using failing stateful computations.
%
We keep a state threaded through the computation, accumulating information to
generate witnesses of fraud.
%
Also, this computations can fail.
%
Finishing  computation threads succefully lead to fraud-proof token creation.

This module is just prototyped and it not very well documented.
%
We refer the reader to the corresponding DocBuild section in the Lean4 code.

The main file is for the scripting language is
``FMMidgard/ProofProtocol/Catalogue/ScriptDef.lean''.
%
For example on how to use it, check
``FMMidgard/ProofProtocol/Catalogue/Ledger/AllInputsValid.lean''.



\subsection{Validity of FraudProofs}

In the Midgard specification, fraud proofs are presenting by stating the desired
property to check, and then decomposed into simpler smaller properties.
%
For each of these simpler properties, a way to generate a witness proving the
violation of the property, a fraud-proof.

Since we do not have a Cardano-ledger abstraction, nor it is very clear how to
check the validity of these properties.
%
We explore a simple idea trying to verify logically that the above
procedure is correct.
%
To see a full documentation we refer the reader to the docbuild documentation in
the file: ``FMMidgard/ProofProtocol/Catalogue/Ledger/AllInputsValid.lean''.


%%FILE blueprint/src/modules/Scheduler.tex
\section{Scheduler}

The scheduler is the module of Midgard in charge of assigning slots of time so
the operator in charge can submit blocks to the Midgrad State Queue.
%
Operators are chosen from the \emph{active queue} from the Operators Directory
following a \emph{round robin} strategy.
%
The active operator queue is sorted by their key.

When we reach the end of the active operators queue, the scheduler is
\emph{rewound}.
%
Therefore the scheduler supports two operations:
\begin{description}
  \item[Advance:] advances the scheduler.
  \item[Rewind:] rewinds the scheduler.
  \end{description}

The Midgard scheduler specification was incomplete, there was no initial state
for the scheduler set.

\subsection{Plots}

\begin{figure}
  \caption{Scheduler Finite State Machine.}
  \includegraphics[width=\textwidth]{Scheduler}
  \label{fig:scheduler:fsm}
\end{figure}

\begin{figure}
  \caption{Scheduler IO plot.}
  \includegraphics[width=\textwidth]{io.scheduler}
  \label{fig:opdir}
\end{figure}

\subsection{Formalizaton}

As state, the scheduler keeps track of the current operator's id plus a posix
time indicating when the current shift starts.
%
The duration of shifts is a protocol parameter.

\begin{definition}[Scheduler State]
  %
  \lean{Scheduler.State.SchedulerDatum}
  \label{dfn:Scheduler.State.SchedulerDatum}
  \leanok
  %
  The scheduler state is simply a record of values:
  \begin{description}
      \item[operator\_id] indicating the current operator's id
      \item[shift\_start] indicating when the shift starts.
    \end{description}
  \end{definition}


\begin{definition}[State Transition Function]
  \lean{Scheduler.StateMachine.next}
  \leanok
  \label{dfn:Scheduler.StateMachine.next}
  %%
  State machine producing the next state triggered by two operations:
  \begin{enumerate}
      \item Advance
      \item Rewind
    \end{enumerate}
  \end{definition}

\begin{definition}[Inductive Relation]
  \lean{Scheduler.Relation}
  \leanok
  \label{dfn:Scheduler.Relation}
  %%
  Similar to the state transition function.
  \end{definition}

\begin{theorem}[Computational Ind. Relation]
  \lean{Scheduler.Computational.computational}
  \uses{dfn:Scheduler.StateMachine.next, dfn:Scheduler.Relation}
  \leanok
  \label{thm:Scheduler.Computational.computational}
  %%
  The state transition function realizes the inductive relation.
  \end{theorem}

\subsection{Properties}

\begin{theorem}[Shifts progresses by shift duration ]
  \leanok
  \lean{Scheduler.Properties.ShiftStepShiftInterval}
  \uses{dfn:MidgardParameters,dfn:InputInfo,dfn:OperatorDirectory.State.OperatorDirectory, dfn:Scheduler.Relation}
  \label{thm:Scheduler.Properties.ShiftStepShiftInterval}
  %%
  At each step, the new shift starts when the previous ended shifted by
\emph{shift\_duration}.
  \end{theorem}

\begin{lemma}[Current Operators was Active]
  \leanok
  \lean{Scheduler.Properties.CurrentOperatorIsChoosenFromActive}
  \uses{dfn:MidgardParameters,dfn:InputInfo,dfn:OperatorDirectory.State.OperatorDirectory, dfn:Scheduler.Relation}
  \label{lmm:Scheduler.Properties.CurrentOperatorIsChoosenFromActive}
  %%
  The scheduler chooses operators from the active queue in the operator directory.
  \end{lemma}

To prove that the current operating operator is active, we need to compose the
scheduler with other modules.
%
In fact, the current operator can be deactivated leading to inconsistent scheduler states.

\subsubsection{Scheduler Progress Theorems}

To prove that we can always make the scheduler tick, we need an additional
condition: there is an active operator.
%
In this case, this condition is a bit more complex, because we may need to rewind.

Rewind is enabled if there are no eligible registered operators waiting to be
activated.
%
There may be eligible operators depending when we are looking at the operators
directory, and the time is defined by the (L1-)transaction's internal validity
upper bound.
%
The specification is not clear about what the criteria to activate registered
operator is.
%
We submitted an issue regarding this.

There are two ways of showing progress in the scheduler.
%
The next active operator in line can request control of the scheduler, and if
the \emph{time has passed}, another operator can make the scheduler advance.

\begin{theorem}[Self Claim Scheduler without Rewinding]
  \lean{Scheduler.Properties.ClaimSchedulerNotRewind}
  \leanok
  \uses{dfn:Scheduler.State.SchedulerDatum,dfn:ComputationalRelation.StateMachine.FromInit%
    ,dfn:ComputationalRelation.StateMachine.AfterInit%
  }
  \label{thm:Scheduler.Properties.ClaimScheduler}
  %
  Assuming we can build a transaction \(\mathsf{tx}\) such that the time
validity bounds in \(\mathsf{tx}\) fit inside the interval
\([ \mathsf{sch.shift\_start} + \mathsf{shift\_duration} , \mathsf{sch.shift\_start} + 2*\mathsf{shift\_duration})\),
then the next operator in line can claim control of the scheduler.
  \end{theorem}

The above theorem makes sense if \(\mathsf{shift\_duration}\) is grater than 1.

\begin{theorem}[Scheduler Rewinds]
  \lean{Scheduler.Properties.ClaimSchedulerRewind}
  \leanok
  \label{thm:Scheduler.Properties.ClaimSchedulerRewind}
  %
  The last active operator rewinds the scheduler when there are no registered
operators waiting to be activated.
  \end{theorem}

Where to check there are not operators waiting to be activated, we just check
the last registered operators activation time.


Upon trying to prove the scheduler can always progress or it is always enabled,
we found that there are conditions leading to locking the scheduler.
%
We submitted an issue and described the situation to Anastasia Labs developers.

Since the specification was not complete, we stop our formal verification efforts here.
%
In the Lean4 file stating properties about the scheduler, there are some ideas
on how to keep going with the formalization efforts.

% %
% Either by the agent requesting the scheduler control or other agents making the
% scheduler advance.

% \begin{theorem}[Scheduler Progress]
%   \lean{Scheduler.Properties.LiveActiveOp}
%   \leanok
%   \uses{dfn:Scheduler.State,dfn:ComputationalRelation.StateMachine.FromInit%
%     ,dfn:ComputationalRelation.StateMachine.AfterInit
%     ,dfn:OperatorDirectory.State.OperatorDirectory.nextActiveId?%
%   }
%   \label{thm:Scheduler.Properties.LiveActiveOp}
%   %
%   If the operators active queue is not empty and there are no eligible
% registered operators waiting for activation, then the scheduler progresses.
%   \end{theorem}

% \subsubsection{Scheduler Stuck Conditions}

% If we need to rewind the scheduler and there are eligible registered operators
% for activation, then the scheduler is stuck until we activate \textbf{all}
% eligible registered operators.

% \begin{definition}[Registered Operators Eligible for activation]
%   \lean{Scheduler.Properties.OpsEligibleForAct}
%   \label{dft:Scheduler.Properties.OpsEligibleForAct}
%   \uses{dfn:OperatorDirectory.State.OperatorDirectory}
%   \leanok
%   %%
%   Given a posix time timepstamp \(\mathsf{time}\),
%   we define the registered operators eligible for activation as those in the registered operators queue whose time is grater or equal than \(\mathsf{time}\).
%   \end{definition}

% \martin{Eligible for activation is not well defined.
%     The spec mentions that looking at the last node is enough, but it is not.
%     Submitted an issue but still I am not sure how to proceed.}

% \begin{theorem}[Scheduler Stuck]
%   % \lean{Scheduler.Properties.NoRewind}
%   \uses{dfn:MidgardParameters,dfn:InputInfo%
%         ,dfn:Scheduler.State,dfn:Scheduler.Relation,dft:Scheduler.Properties.OpsEligibleForAct}
%   \leanok
%   \label{thm:Scheduler.Properties.NoRewind}
%   %%
%   If we finished a cycle and there are eligible operators waiting to be
% activated, then there is no action that makes the scheduler advance.
%   \end{theorem}

% \begin{proof}
%   \proves{thm:Scheduler.Properties.NoRewind}
%   \leanok
%   We cannot advance the scheduler because we reached the end of a cycle and we
% cannot rewind because we have eligible operators waiting to be activated.
% %
%   Therefore, we cannot do anything from the point of view of the scheduler.
%   \end{proof}


%%FILE blueprint/src/modules/Settlement.tex
\section{Settlement}

The Midgard settlement module is in charge of keeping tract of user events:
deposits, withdrawals and L1-transactions.
%
Operators when pushing blocks into the State Queue also need to take into
account user events and register them properly.
%
Operators not submitting user events are penalized and their bonds slashed.

Originally, this module was implemented as a list, shadowing the State Queue.
%
The specification recently changed, and now, the settlement module only keeps a
collection of settlement, and each settlement is linked to a node in the state
queue.

We provide a mechanization of this module.

\subsection{Plots}

\begin{figure}
  \caption{Settlement Input/Output plots.}
  \includegraphics[width=\textwidth]{io.setts}
  \label{fig:opdir}
\end{figure}


\subsection{Mechanization}

\begin{definition}[Actions]
  \leanok
  \lean{Settlement.State.Actions}
  \label{dfn:Settlement.State.Actions}
  Actions :
  \begin{itemize}
          \item Spawn
          \item Remove
          \item AttachResolutionClaim
          \item DisproveResolutionClaim
          \end{itemize}
  \end{definition}

\begin{definition}[Relation]
  \leanok
  \lean{Settlement.Relation.Rel}
  \label{dfn:Settlement.Relation.Rel}
  \end{definition}

\begin{definition}[Computation]
  \leanok
  \lean{Settlement.Machine.next}
  \label{dfn:Settlement.Machine.next}
  \end{definition}
  


%%FILE blueprint/src/modules/StateQueue.tex
\section{State Queue}

State queue is the Midgard module in charge of keeping track of the L2-Midgard
blocks.
%
The operator assigned at the scheduler appends new blocks into the state queue,
and after waiting the required period (a Midgard parameter), surviving blocks
commit to the history of the L2.

The queue itself is a \emph{heterogeneous} list.
%
The root node contains a summary of the current Midgard blocks.
%
The rest of the list is a queue of blocks waiting to be committed to the
history.

\subsection{Requires}

\paragraph{Confirmed State.}
%
The root node of the State Queue keeps a summary of the history of Midgard.
%
It keeps track of the previous header hash, utxo set root hash, the current
header hash, protocol version and start and end time.

\begin{definition}{Confirmed State}
  \label{dfn:StateQueue.State.ConfirmedState}
  \lean{StateQueue.State.ConfirmedState}
  \leanok
  Confirmed States are records containing the initial starting time of the protocol,
  the last recorded end time (last's confirmed block end time),
  the last confirmed block previous hash,
  the last confirmed block hash,
  the root of the merkle tree representing the set of UTXOs,
  and the protocol version.
\end{definition}

\paragraph{Unordered Simpled Linked Lists}
%
  The queue itself is an unordered simple linked list.
  %
  The heterogeneous list is represented simply as the disjoint union of
  Confirmed States and \emph{Hash header} nodes.

\paragraph{FraudProof tokens}
%
  Additionally, to remove fraudulent blocks, we need to keep track of Fraud
  Proof tokens.

\paragraph{Scheduler}
%
We need to interact with the scheduler to know who the current operator is.

\subsection{Formalization}

\paragraph{State}
%
The state of the State Queue is a Midgard Unordered linked list.
%
We have the additional strong condition that keys are unique.
%
To each block, the key is defined as the hash of the block.
%
Since blocks contain the previous block hash, the resulting hash has very low
chances of collisions.
%
However, the study of hash collisions are left outside this formalization.

\martin{State Transitions}
\begin{definition}[State Transition Actions]
  \label{dfn:StateQueue.State.Actions}
  \lean{StateQueue.State.Actions}\leanok
  Operators can interact with the state queue as follows:
  \begin{description}
    \item[CommitBlock] commiting a new block to the queue
    \item[MergeToConfirmed]  merging the first waiting block into the history
    \item[RemoveFraudulentBlockRetired] removing fraudulent blocks
    \item[RemoveFraudulentBlockActive] removing fraudulent blocks
  \end{description}
\end{definition}

The state queue module interacts tightly with the Operator Directory, and such
interaction is defined when composing the two modules.

%% Place Holder definition.
\begin{definition}[Inductive Relation]
  \label{dfn:StateQueue.Relation}
  \lean{StateQueue.Relation}\leanok
  \uses{dfn:StateQueue.State.Actions, dfn:FraudProofTokens.State.FraudProofTokens, dfn:Scheduler.State.SchedulerDatum}
\end{definition}

\begin{theorem}[State Queue inductive relation is computational]
  \label{thm:StateQueue.Computational.computational}
  \lean{StateQueue.Computational.computational}
  \leanok
  \uses{thm:Scheduler.Computational.computational, thm:OperatorDirectory.Computational.ComputationalRelation}
\end{theorem}

\subsection{Properties}

\begin{theorem}
  \label{thm:StateQueue.Props.only_one_confirmed_state}
  \lean{StateQueue.Props.only_one_confirmed_state}
  \leanok
  \uses{dfn:MUList.State.ListToken, dfn:ComputationalRelation.StateMachine.FromInit}
  The committed history of Midgard is at the root node.
  \end{theorem}

\begin{theorem}
  \label{thm:StateQueue.Props.block_merging_order}
  \lean{StateQueue.Props.block_merging_order}
  \leanok
  The block being merged into the history is the first node in the state queue
list.
  \end{theorem}

If keys are unique, we cannot skip elements.

\subsubsection{Invariants}



%%FILE blueprint/src/modules/UserEvent.tex
\section{User Events}

The Midgrad User Events module is in charge of handling communications between
Cardano-L1 and Midgard-L2.

The specification change since we begun our formalization efforts, and we only
prototyped this module.
%
For a better description we refer the reader to the DocBuild documentation:
``FMMidgard/Bridge/State.lean'' and
``FMMidgard/Bridge/Relation.lean''.

